\chapter[Liebesgaben für die Taubstummenanstalt]{Erhalten an Liebesgaben für die Taubstummenanstalt zu North Detroit, Mich}
\quellseite{15}

(Vom 22. October 1902 bis 1. December 1903.)

Aus Michigan. Detroit: R. R. durch P. Smutal 2 getr., noch gute Mädchenjacken. Von P. Smutal und durch ihn: 5 alte Hüte, 5 M.-Waists, 1 Winterjacke, 1 Rock u. Weste; für Ella und Frida Stelt: 5 Kleider, 1 Hut, 1 Winterjacke, 1 Cape, 1 weißen Unterrock. Chas. Manzelmann 1 Dutz. Besen, 40 Besenstiele, 1 alten Ofen, 4 Reisbesen. Frauenver. für Taubstummenunterstützung 4 M.-Hemden, 3 Pr. Unterhosen, 1 Dutz. Spulen Zwirn, 2 Bat. Nähnadeln, 2 Anzüge, Unterzeug, die ganze Weihnachtsbescherung (Baum und Zubehör, Geschenke für jedes Kind besonders u. für alle im Allgem.), 3 Spiegel, die Tapetirung des Kranken- u. Fremdenzimmers, 1 lawn mower, für arme Kinder \$7.29, 10 galv. Eimer, 3 Blechschüsseln, 2 emaillirte Töpfe, 8 Näpfe, 2 Bratpfannen, 2 Dutz. Eßlöffel, Confirmationskleider für 2 Mädchen. A. Stendel 1 Laterne, 7 dust pans, 3 pokers. Fr. Otto einige getr. Kleider u. Schuhe. Frau R. R. 1 einger. Bild u. einige getr. Kleider. Fr. J. C. 1 gebr. Hängeleuchte. Fr. Kieke Blumenarbeit für 4.00. P. Sagens Frauenver. einige comforters, Wolldecken, Schürzen, Knabenblusen; für J. Klatt 3 Pr. Schuhe, 2 Blousen, 1 Anzug. C. Wendt 2 Pr. Schuhe, 1 Pr. Ueberschuhe. Fr. Taube 1 Schinken. Frau Keyl 36 Handspiegel. Frau Flemming 6 Kissenbezüge. Mehrere Frauen in P. Ottes Gem. 4 Paar Kleider. J. Reindel getr. Unterzeug, einige Hemden, Manschetten, Kragen, Schleifen. Fr. Knorr 3 Flaschen Hien Fong.

— Frau D. Frank, Cottonwood, Minn., 5.00. Frauenver. Steiner, Mich., 3 comforters. Frau G. Heiß u. A. Reingel dasel. je 1 comforter. Frau Geo. Fiedler u. G. Reinhardt dasel. je 1 Wolldecke. G. Zorr, Toledo, u. Hr. Debur, Detroit, je 1.00; J. Gwerth, Akron, O., 3.00; J. G. Dahms, Sibley, Iowa, .50; J. Felske, Augsburg, Ont., .25 (alles Geld bestimmt für die Weihnachtsbescherung). Thoms Bros., Chicago, für Weihnachten 1 Groß Stahlfedern, 1 Groß Faber-Bleistifte, 3 Dutz. Notizbücher, 7 Dutz. scratch blocks, 3 Dutz. feine Federhalter. Frau Ch. Knöll u. Frau Leo Durandt, Chicago, 1.00, resp. 2.00. H. Meyer, Pigeon, Mich., 2 Gänse, 15 Pfd. Kalbfleisch. J. Gwerth, Akron, O., eine gute Hängeleuchte. Schult. in Pigeon, Mich., 16.80, Kostgeld für Anna Meier. R. Stolzenfeld, North Detroit, 1½ Bu. Saatkorn. Von der Gem. das. Fahren des Feuerholzes.
\quellseite{16}
A. Paul, Buffalo, 1.00. Frl. M. Tönsing, Cleveland, O., 5.00. Frau W. Rasmus, Bloomer, Wis., 4.00. Von den Frauen N. N. u. N. N., Balmya, Mo., je 1.00. N. N., Cleveland, O., 5.00. Herr u. Frau Renner, Frau Knößl, Chicago, je 1.00. Hr. Leisheit, Yorkville, Ill., 1.00. G. Engel, North Detroit, 3 Pfd. Honig, 1 Pfd. kleine Zwiebeln, ½ Bu. grüne Zwiebeln, 6 Pfd. Wurst. Gera-Nährkränzchen durch Lehrer Seile, Gera, Mich., 5.00 für Heizungsanlage. Gem. in Farmington, Mo., 15 Bu. Kart. An Zeitschriften erhalten: Lutheraner, Missions-Taube, Kinder- und Jugendblatt, Pioneer, Lutheran Witness, für die Kleinen, Gemeindeblatt, das Blatt der Michigan-Kinderfreundgesellschaft, Die Abendschule, Illustrated Home Journal, Rundschau, Germania, Haus- u. Bauernfreund, Lutherische Kirchenzeitung, Luth. Sunday School Herald, Lutheran Guide, The Little Ones, Junior S.-S. Lesson Book, Augsburg S.-S. Teacher, Michigan Mirror, Kentucky Standard, Illinois’ New Era, The Deaf Mute Voice (Jackson, Mich.), Detroit Abendpost, Detroit Journal. Gott wolle es allen lieben Gebern reichlich vergelten.

W. Gielow.
